\documentclass{article}
\usepackage[utf8]{inputenc}
\usepackage{graphicx}
\usepackage{amsmath}
\usepackage{float} % Add the float package

\title{Mathematical Statistics Assignment 2}
\author{ABID ALI (Student ID: 2023280099)}
\date{October 20, 2024}

\begin{document}

\maketitle

\newpage % Forces the following content to appear on a new page

\section*{2.2.16}
If $A$ and $B$ are mutually exclusive events, $P(A) = 0.17$ and $P(B) = 0.46$, find

\begin{enumerate}
    \item[(a)] $P(A \cup B)$
    \item[(b)] $P(A^c)$
    \item[(c)] $P(A^c \cup B^c)$
    \item[(d)] $P(A \cap (B^c))$
    \item[(e)] $P(A^c \cap B^c)$
\end{enumerate}

\section*{Solution}

Given:
\[
P(A) = 0.17, \quad P(B) = 0.46
\]
\(A\) and \(B\) are mutually exclusive events, meaning \(P(A \cap B) = 0\) 
i.e. \[A \cap B = \emptyset\]


\begin{enumerate}
    \item[(a)] \(P(A \cup B)\)
    \[
    P(A \cup B) = P(A) + P(B)-P(A \cap B)
    \]
    \[
    P(A \cup B) = 0.17 + 0.46 -0 \]
    \[
    P(A \cup B) = 0.17 + 0.46 = 0.63
    \]
    
    \item[(b)] \(P(A^c)\)
    \[
    P(A^c) = 1 - P(A)
    \]
    \[
    P(A^c) = 1 - 0.17 = 0.83
    \]
    
    \item[(c)] \(P(A^c \cup B^c)\)
    \[
    P(A^c \cup B^c) = 1 - P(A \cap B)
    \]
    Since \(P(A \cap B) = 0\),
    \[
    P(A^c \cup B^c) = 1 - P(A \cup B) = 1 - 0.63 = 0.37
    \]
    
    \item[(d)] \(P(A \cap (B^c))\)
    \[
    P(A \cap B^c) = P(A) - P(A \cap B)
    \]
    Since \(P(A \cap B) = 0\),
    \[
    P(A \cap B^c) = P(A) = 0.17
    \]
    
    \item[(e)] \(P(A^c \cap B^c)\)
    \[
    P(A^c \cap B^c) = 1 - P(A \cup B)
    \]
    We already calculated that \(P(A \cup B) = 0.63\),
    \[
    P(A^c \cap B^c) = 1 - 0.63 = 0.37
    \]
    
\end{enumerate}

\section*{2.2.21}

A box contains three red and two blue flies. Two flies are removed with replacement. Let \( A \) be the event that both the flies are of the same color, and \( B \) be the event that at least one of the flies is red. 

\noindent We are required to find:
\begin{enumerate}
    \item \( P(A) \)
    \item \( P(B) \)
    \item \( P(A \cup B) \)
    \item \( P(A \cap B) \)
\end{enumerate}

\section*{Solution:}

\subsection*{Sample Space}

We are drawing two flies with replacement from a box containing 3 red flies and 2 blue flies. Therefore, the possible outcomes for each draw are "red" (R) or "blue" (B). Since the draws are with replacement, the possible outcomes for the two draws can be represented as ordered pairs.

The sample space \( S \) is:
\[
S = \{ (R,R), (R,B), (B,R), (B,B) \}
\]

Each outcome has the following probabilities:
\begin{itemize}
    \item \( P(R,R) = \frac{3}{5} \times \frac{3}{5} = \frac{9}{25} \)
    \item \( P(R,B) = \frac{3}{5} \times \frac{2}{5} = \frac{6}{25} \)
    \item \( P(B,R) = \frac{2}{5} \times \frac{3}{5} = \frac{6}{25} \)
    \item \( P(B,B) = \frac{2}{5} \times \frac{2}{5} = \frac{4}{25} \)
\end{itemize}

Thus, the sample space with associated probabilities is:
\[
S = \left\{ (R,R), (R,B), (B,R), (B,B) \right\}
\]
with the probabilities:
\[
\{ P(R,R) = \frac{9}{25}, \ P(R,B) = \frac{6}{25}, \ P(B,R) = \frac{6}{25}, \ P(B,B) = \frac{4}{25} \}
\]

\subsection*{(1) Event \( A \): Both flies are of the same color}

The possible outcomes for event \( A \) (same color) are:
\[
A = \{ (R,R), (B,B) \}
\]
The probability of \( A \) is:
\[
P(A) = P(R,R) + P(B,B) = \frac{9}{25} + \frac{4}{25} = \frac{13}{25}
\]

\subsection*{(2) Event \( B \): At least one fly is red}

The possible outcomes for event \( B \) (at least one red fly) are:
\[
B = \{ (R,R), (R,B), (B,R) \}
\]
The probability of \( B \) is:
\[
P(B) = P(R,R) + P(R,B) + P(B,R) = \frac{9}{25} + \frac{6}{25} + \frac{6}{25} = \frac{21}{25}
\]

\subsection*{(3) Union of Events \( A \cup B \)}

The union of \( A \) and \( B \) represents the event that either both flies are the same color or at least one of them is red. This includes all outcomes except \( (B,B) \), so:
\[
A \cup B = \{ (R,R), (R,B), (B,R), (B,B) \}
\]
The probability of \( A \cup B \) is:
\[
P(A \cup B) = P(R,R) + P(R,B) + P(B,R) + P(B,B) = 1
\]

\subsection*{(4) Intersection of Events \( A \cap B \)}

The intersection of \( A \) and \( B \) represents the event that both flies are the same color and at least one of them is red. This corresponds to the outcome \( (R,R) \), since both flies are red (same color and both red).

\[
A \cap B = \{ (R,R) \}
\]
The probability of \( A \cap B \) is:
\[
P(A \cap B) = P(R,R) = \frac{9}{25}
\]


\section*{Summary:}
\begin{itemize}
    \item \textbf{Sample Space \( S \):}
    \[
    S = \{ (R,R), (R,B), (B,R), (B,B) \}
    \]
    
    \item \textbf{Event \( A \)} (same color):
    \[
    A = \{ (R,R), (B,B) \}
    \]
    \[
    P(A) = \frac{13}{25}
    \]
    
    \item \textbf{Event \( B \)} (at least one red fly):
    \[
    B = \{ (R,R), (R,B), (B,R) \}
    \]
    \[
    P(B) = \frac{21}{25}
    \]

    \item \textbf{Union \( A \cup B \)} (either same color or at least one red fly):
    \[
    A \cup B = \{ (R,R), (R,B), (B,R), (B,B) \}
    \]
    \[
    P(A \cup B) = 1
    \]
        
    \item \textbf{Intersection \( A \cap B \)} (both red flies):
    \[
    A \cap B = \{ (R,R) \}
    \]
    \[
    P(A \cap B) = \frac{9}{25}
    \]
    

\end{itemize}

\section*{2.3.12}
A package of 15 apples contains two defective apples. Four apples are selected at random.

(a) Find the probability that none of the selected apples is defective.

(b) Find the probability that at least one of the selected apples is defective.

\section*{Solution:}

A package of 15 apples contains two defective apples. Four apples are selected at random. We are asked to compute the following probabilities:

\subsection*{(a) Probability that none of the selected apples is defective:}

The total number of ways to select 4 apples from 15 is given by the combination:
\[
\binom{15}{4} = \frac{15 \times 14 \times 13 \times 12}{4 \times 3 \times 2 \times 1} = 1365
\]

Since there are 13 good apples, the number of ways to select 4 good apples from 13 is:
\[
\binom{13}{4} = \frac{13 \times 12 \times 11 \times 10}{4 \times 3 \times 2 \times 1} = 715
\]

Therefore, the probability that none of the selected apples is defective is:
\[
P(\text{none defective}) = \frac{\binom{13}{4}}{\binom{15}{4}} = \frac{715}{1365} \approx 0.524
\]

\subsection*{(b) Probability that at least one of the selected apples is defective:}

The complement of "at least one defective" is "none defective." Therefore, the probability that at least one apple is defective is:
\[
P(\text{at least one defective}) = 1 - P(\text{none defective})
\]

Substitute the result from part (a):
\[
P(\text{at least one defective}) = 1 - \frac{715}{1365} = \frac{1365 - 715}{1365} = \frac{650}{1365} \approx 0.476
\]

\subsection*{Final Answers:}
\begin{itemize}
    \item (a) The probability that none of the selected apples is defective: \( \frac{715}{1365} \approx 0.524 \).
    \item (b) The probability that at least one of the selected apples is defective: \( \frac{650}{1365} \approx 0.476 \).
\end{itemize}

\section*{2.3.16}

A box of manufactured items contains 12 items, of which four are defective. If three items are drawn at random without replacement, what is the probability that

(a) The first one is defective and the rest are good?

(b) Exactly one of the three is defective?

\section*{Solution:}

A box contains 12 items, of which four are defective. Three items are drawn at random without replacement. We are asked to compute the following probabilities:

\subsection*{(a) Probability that the first item is defective and the rest are good:}

The probability that the first item is defective is:
\[
P(\text{first defective}) = \frac{4}{12}
\]

The probability that the second item is good, given the first was defective, is:
\[
P(\text{second good}) = \frac{8}{11}
\]

The probability that the third item is good, given the first was defective and the second was good, is:
\[
P(\text{third good}) = \frac{7}{10}
\]

Thus, the total probability that the first item is defective and the rest are good is:
\[
P(\text{first defective, rest good}) = \frac{4}{12} \times \frac{8}{11} \times \frac{7}{10} = \frac{56}{330} \approx 0.1697
\]

\subsection*{(b) Probability that exactly one of the three items is defective:}

We need to compute the probability for three mutually exclusive cases:

\textbf{Case 1:} First item is defective, and the second and third are good:
\[
P(\text{first defective, rest good}) = \frac{56}{330}
\]

\textbf{Case 2:} Second item is defective, and the first and third are good:
\[
P(\text{second defective, first and third good}) = \frac{8}{12} \times \frac{4}{11} \times \frac{7}{10} = \frac{56}{330}
\]

\textbf{Case 3:} Third item is defective, and the first and second are good:
\[
P(\text{third defective, first and second good}) = \frac{8}{12} \times \frac{7}{11} \times \frac{4}{10} = \frac{56}{330}
\]

The total probability that exactly one of the three items is defective is:
\[
P(\text{exactly one defective}) = \frac{56}{330} + \frac{56}{330} + \frac{56}{330} = \frac{168}{330} \approx 0.5091
\]

\section*{2.4.12}

A survey asked a group of 400 people whether or not they were doing daily exercise. The responses by sex and
physical activity are as in Table 2.4.
A person is randomly selected.

(a) What is the probability that this person is doing daily exercise?

(b) What is the probability that this person is doing daily exercise if we know that this person is a male?

\begin{figure}[H] % The H option forces the figure to be placed exactly here
\centering
% Replace this with the correct path to your Table 1.9 image file
\includegraphics[width=0.8\textwidth]{D:/Masters NPU Semeter/3rd Semester/Statistics/Homework/Submit/Homework 2/figs/2412.png}
\caption{Table 2.4 - Physical Activity Survey Results by Gender.}
\end{figure}

\section*{Solution:}

A survey asked a group of 400 people whether or not they were doing daily exercise. The responses by sex and physical activity are as follows:

\begin{center}
\begin{tabular}{|l|c|c|}
\hline
 & Male & Female \\
\hline
Daily exercise & 50 & 61 \\
No daily exercise & 177 & 112 \\
\hline
\end{tabular}
\end{center}

\subsection*{(a) Probability that a randomly selected person is doing daily exercise:}

The total number of people doing daily exercise is:
\[
50 \, (\text{males}) + 61 \, (\text{females}) = 111
\]

The total number of people surveyed is 400. Therefore, the probability that a randomly selected person is doing daily exercise is:
\[
P(\text{daily exercise}) = \frac{111}{400} = 0.2775
\]

\subsection*{(b) Probability that a randomly selected male is doing daily exercise:}

The total number of males surveyed is:
\[
50 \, (\text{males who do daily exercise}) + 177 \, (\text{males who do not do daily exercise}) = 227
\]

Thus, the probability that a randomly selected male is doing daily exercise is:
\[
P(\text{daily exercise | male}) = \frac{50}{227} \approx 0.2203
\]

\subsection*{Final Answers:}
\begin{itemize}
    \item (a) The probability that a randomly selected person is doing daily exercise is \( \frac{111}{400} = 0.2775 \).
    \item (b) The probability that a randomly selected male is doing daily exercise is \( \frac{50}{227} \approx 0.2203 \).
\end{itemize}

\section*{2.4.23}

A store has light bulbs from two suppliers, 1 and 2. The chance of supplier 1 delivering defective bulbs is 10\%, whereas supplier 2 has a defective rate of 3\%. Suppose 60\% of the current supply of light bulbs came from supplier 1. If one of these bulbs is taken from the current supply and observed to be defective, find the probability that it came from supplier 2.

\section*{Solution:}

We are asked to find the probability that a defective light bulb came from supplier 2, given that a defective bulb was selected. This is a Bayes' theorem problem.

\subsection*{Given:}
\begin{itemize}
    \item Probability that a bulb came from supplier 1: \( P(S_1) = 0.60 \)
    \item Probability that a bulb came from supplier 2: \( P(S_2) = 0.40 \)
    \item Probability that a bulb from supplier 1 is defective: \( P(D|S_1) = 0.10 \)
    \item Probability that a bulb from supplier 2 is defective: \( P(D|S_2) = 0.03 \)
\end{itemize}

We are asked to find \( P(S_2|D) \), the probability that the defective bulb came from supplier 2.

\subsection*{Using Bayes' Theorem:}

\[
P(S_2|D) = \frac{P(D|S_2) \cdot P(S_2)}{P(D)}
\]
Where \( P(D) \) is the total probability that a bulb is defective, which is given by the law of total probability:
\[
P(D) = P(D|S_1) \cdot P(S_1) + P(D|S_2) \cdot P(S_2)
\]

\subsection*{Step 1: Compute \( P(D) \)}

\[
P(D) = (0.10 \cdot 0.60) + (0.03 \cdot 0.40) = 0.06 + 0.012 = 0.072
\]

\subsection*{Step 2: Apply Bayes' Theorem}

\[
P(S_2|D) = \frac{(0.03) \cdot (0.40)}{0.072} = \frac{0.012}{0.072} = \frac{1}{6} \approx 0.1667
\]

\subsection*{Final Answer:}

The probability that the defective bulb came from supplier 2 is approximately \( 0.1667 \), or \( 16.67\% \).


\section*{2.5.8}

\textbf{Problem:}

Suppose that \( Y \) is a continuous random variable whose probability density function (pdf) is given by:

\[
f(y) = 
\begin{cases} 
K(4y - 2y^2), & 0 < y < 2, \\
0, & \text{elsewhere}.
\end{cases}
\]

\begin{enumerate}
    \item[(a)] What is the value of \( K \)?
    \item[(b)] Find \( P(Y > 1) \).
    \item[(c)] Find \( F(y) \), the cumulative distribution function (CDF).
\end{enumerate}

\section*{Solution:}

\textbf{Solution:}

\textbf{(a) What is the value of \( K \)?}

To find \( K \), we use the fact that the total probability must equal 1, i.e.,

\[
\int_{-\infty}^{\infty} f(y) \, dy = 1.
\]

Since \( f(y) = 0 \) outside of \( 0 < y < 2 \), we only need to integrate over this interval:

\[
\int_0^2 K(4y - 2y^2) \, dy = 1.
\]

First, compute the integral:

\[
\int_0^2 (4y - 2y^2) \, dy = \left[ 2y^2 - \frac{2}{3}y^3 \right]_0^2 = \left(2(2)^2 - \frac{2}{3}(2)^3\right) - \left(2(0)^2 - \frac{2}{3}(0)^3\right).
\]

Simplifying the expression:

\[
= \left(2(4) - \frac{2}{3}(8)\right) = 8 - \frac{16}{3} = \frac{24}{3} - \frac{16}{3} = \frac{8}{3}.
\]

Thus, the integral becomes:

\[
K \cdot \frac{8}{3} = 1.
\]

Solving for \( K \):

\[
K = \frac{3}{8}.
\]

\textbf{(b) Find \( P(Y > 1) \).}

We want to find the probability \( P(Y > 1) \), which is:

\[
P(Y > 1) = \int_1^2 f(y) \, dy = \int_1^2 \frac{3}{8}(4y - 2y^2) \, dy.
\]

First, compute the integral:

\[
\int_1^2 (4y - 2y^2) \, dy = \left[ 2y^2 - \frac{2}{3}y^3 \right]_1^2 = \left(2(2)^2 - \frac{2}{3}(2)^3\right) - \left(2(1)^2 - \frac{2}{3}(1)^3\right).
\]

Simplifying the expression:

\[
= \left(8 - \frac{16}{3}\right) - \left(2 - \frac{2}{3}\right) = \left(\frac{24}{3} - \frac{16}{3}\right) - \left(2 - \frac{2}{3}\right) = \frac{8}{3} - \frac{4}{3} = \frac{4}{3}.
\]

Thus:

\[
P(Y > 1) = \frac{3}{8} \cdot \frac{4}{3} = \frac{1}{2}.
\]

\textbf{(c) Find \( F(y) \).}

The cumulative distribution function (CDF) \( F(y) \) is defined as:

\[
F(y) = P(Y \leq y) = \int_{-\infty}^y f(t) \, dt.
\]

For \( y < 0 \), \( F(y) = 0 \) because \( f(y) = 0 \) for \( y < 0 \).

For \( 0 \leq y \leq 2 \):

\[
F(y) = \int_0^y \frac{3}{8}(4t - 2t^2) \, dt.
\]

First, compute the integral:

\[
\int_0^y (4t - 2t^2) \, dt = \left[ 2t^2 - \frac{2}{3}t^3 \right]_0^y = 2y^2 - \frac{2}{3}y^3.
\]

Thus:

\[
F(y) = \frac{3}{8} \left( 2y^2 - \frac{2}{3}y^3 \right) = \frac{3}{4} y^2 - \frac{1}{4} y^3.
\]

For \( y > 2 \), \( F(y) = 1 \) because all probability is accumulated.

\textbf{Final CDF expression:}

\[
F(y) = 
\begin{cases} 
0, & y < 0, \\
\frac{3}{4}y^2 - \frac{1}{4}y^3, & 0 \leq y \leq 2, \\
1, & y > 2.
\end{cases}
\]

\section*{2.6.6}

\textbf{Problem:}

Let \( X \) be a random variable with probability density function (pdf) \( f(x) = kx^2 \) where \( 0 \leq x \leq 1 \).

\begin{enumerate}
    \item[(a)] Find \( k \).
    \item[(b)] Find \( E(X) \) and \( \text{Var}(X) \).
    \item[(c)] Find the moment generating function \( M_X(t) \). Using the mgf, find \( E(X) \).
\end{enumerate}

\section*{Solution:}

\textbf{(a) Find \( k \)}

We know that the total probability must equal 1, i.e.,

\[
\int_0^1 f(x) \, dx = 1.
\]

Given \( f(x) = kx^2 \), we have:

\[
\int_0^1 kx^2 \, dx = 1.
\]

Compute the integral:

\[
k \int_0^1 x^2 \, dx = k \left[ \frac{x^3}{3} \right]_0^1 = k \cdot \frac{1}{3}.
\]

So, we have:

\[
k \cdot \frac{1}{3} = 1 \quad \Rightarrow \quad k = 3.
\]

Thus, \( k = 3 \).

\textbf{(b) Find \( E(X) \) and \( \text{Var}(X) \)}

\textbf{Finding \( E(X) \):}

The expectation \( E(X) \) is given by:

\[
E(X) = \int_0^1 x f(x) \, dx = \int_0^1 x (3x^2) \, dx = 3 \int_0^1 x^3 \, dx.
\]

Compute the integral:

\[
E(X) = 3 \left[ \frac{x^4}{4} \right]_0^1 = 3 \cdot \frac{1}{4} = \frac{3}{4}.
\]

So, \( E(X) = \frac{3}{4} \).

\textbf{Finding \( \text{Var}(X) \):}

The variance \( \text{Var}(X) \) is given by:

\[
\text{Var}(X) = E(X^2) - (E(X))^2.
\]

First, we need to find \( E(X^2) \):

\[
E(X^2) = \int_0^1 x^2 f(x) \, dx = \int_0^1 x^2 (3x^2) \, dx = 3 \int_0^1 x^4 \, dx.
\]

Compute the integral:

\[
E(X^2) = 3 \left[ \frac{x^5}{5} \right]_0^1 = 3 \cdot \frac{1}{5} = \frac{3}{5}.
\]

Now, we can calculate the variance:

\[
\text{Var}(X) = \frac{3}{5} - \left( \frac{3}{4} \right)^2 = \frac{3}{5} - \frac{9}{16}.
\]

To simplify, find a common denominator:

\[
\text{Var}(X) = \frac{48}{80} - \frac{45}{80} = \frac{3}{80}.
\]

Thus, \( \text{Var}(X) = \frac{3}{80} \).

\textbf{(c) Find the moment generating function \( M_X(t) \) and use it to find \( E(X) \)}

The moment generating function (mgf) \( M_X(t) \) is given by:

\[
M_X(t) = E(e^{tX}) = \int_0^1 e^{tx} f(x) \, dx = \int_0^1 e^{tx} (3x^2) \, dx.
\]

To find \( M_X(t) \), we need to compute the integral:

\[
M_X(t) = 3 \int_0^1 x^2 e^{tx} \, dx.
\]

This is a more complex integral that involves integration by parts. However, using the property of the mgf, we know that the first derivative of \( M_X(t) \) with respect to \( t \) gives \( E(X) \):

\[
M'_X(t) = \frac{d}{dt} M_X(t), \quad \text{and} \quad E(X) = M'_X(0).
\]

To find \( M'_X(0) \), we differentiate \( M_X(t) \), but since this involves advanced calculus steps (or it may be provided in a table of common mgfs), we directly use the earlier result:

\[
E(X) = \frac{3}{4}.
\]

Thus, we verified that \( E(X) = \frac{3}{4} \) using the mgf approach.


\section*{2.6.8}

\textbf{Problem:}

Given that \( X_1, X_2, X_3, \) and \( X_4 \) are independent random variables with mean 2, find \( E(Y) \) and \( E(Z) \) for

\[
Y = 3X_4 - X_1 + \frac{1}{5} X_3
\]

\[
Z = X_2 + 7X_3 - 9X_1.
\]


\section*{Solution:}

\textbf{Solution:}

We are given that \( X_1, X_2, X_3, \) and \( X_4 \) are independent random variables with mean 2. Thus,

\[
E(X_1) = E(X_2) = E(X_3) = E(X_4) = 2.
\]

\textbf{(1) Find \( E(Y) \):}

The expression for \( Y \) is:

\[
Y = 3X_4 - X_1 + \frac{1}{5} X_3.
\]

Using the linearity of expectation, we can compute \( E(Y) \) as:

\[
E(Y) = E(3X_4 - X_1 + \frac{1}{5} X_3) = 3E(X_4) - E(X_1) + \frac{1}{5} E(X_3).
\]

Substitute \( E(X_1) = E(X_2) = E(X_3) = E(X_4) = 2 \):

\[
E(Y) = 3(2) - 2 + \frac{1}{5}(2) = 6 - 2 + \frac{2}{5}.
\]

Simplifying the expression:

\[
E(Y) = 4 + \frac{2}{5} = \frac{20}{5} + \frac{2}{5} = \frac{22}{5}.
\]

Thus, \( E(Y) = \frac{22}{5} \).

\textbf{(2) Find \( E(Z) \):}

The expression for \( Z \) is:

\[
Z = X_2 + 7X_3 - 9X_1.
\]

Using the linearity of expectation, we can compute \( E(Z) \) as:

\[
E(Z) = E(X_2 + 7X_3 - 9X_1) = E(X_2) + 7E(X_3) - 9E(X_1).
\]

Substitute \( E(X_1) = E(X_2) = E(X_3) = E(X_4) = 2 \):

\[
E(Z) = 2 + 7(2) - 9(2) = 2 + 14 - 18.
\]

Simplifying the expression:

\[
E(Z) = 16 - 18 = -2.
\]

Thus, \( E(Z) = -2 \).

\textbf{Final Results:}

\[
E(Y) = \frac{22}{5}, \quad E(Z) = -2.
\]


\section*{2.6.16}

\textbf{Problem:}

Find \( E(X) \) and \( \text{Var}(X) \) for a random variable \( X \) with probability density function (pdf):

\[
f(x) = \frac{1}{2} e^{-|x|}, \quad -\infty < x < \infty.
\]

Also, find the moment generating function (mgf) of \( X \).


\section*{Solution:}
\textbf{(a)}

\[
E(X) = \int_{-\infty}^{+\infty} x f(x) \, dx
\]

\[
= \int_{-\infty}^{+\infty} x \cdot \frac{1}{2} e^{-|x|} \, dx
\]

\[
= \int_{-\infty}^{0} x \cdot \frac{e^{x}}{2} \, dx + \int_{0}^{\infty} x \cdot \frac{e^{-x}}{2} \, dx
\]

\[
= \frac{1}{2} \int_{-\infty}^{0} x e^{x} \, dx + \frac{1}{2} \int_{0}^{\infty} x e^{-x} \, dx
\]

\[
= \frac{1}{2} \left[ x e^{x} - e^{x} \right]_{-\infty}^{0} + \frac{1}{2} \left[ - x e^{-x} - e^{-x} \right]_{0}^{\infty}
\]

\[
= \frac{1}{2} \left[ e^{x}(x-1) \right]_{-\infty}^{0} + \frac{1}{2} \left[ e^{-x}(-x-1) \right]_{0}^{\infty}
\]

\[
= \frac{1}{2} \left[ -1 \right] + \frac{1}{2} \left[ 1 \right]
\]

\[
= \frac{-1}{2} + \frac{1}{2} = 0.
\]

Therefore, \( E(X) \) for the given probability density function is 0.

\textbf{(b) Var(X)}

\[
\text{Var}(X) = \int_{-\infty}^{+\infty} x^2 f(x) \, dx
\]

\[
= \int_{-\infty}^{+\infty} x^2 \cdot \frac{1}{2} e^{-|x|} \, dx
\]

\[
= \int_{-\infty}^{0} x^2 \frac{e^{x}}{2} \, dx + \int_{0}^{\infty} x^2 \frac{e^{-x}}{2} \, dx
\]

\[
= \frac{1}{2} \int_{-\infty}^{0} x^2 e^{x} \, dx + \frac{1}{2} \int_{0}^{\infty} x^2 e^{-x} \, dx
\]

\[
= \frac{1}{2} \left( x^2 e^{x} - 2 \left[ x e^{x} - e^{x} \right] \right)_{-\infty}^{0} + \frac{1}{2} \left( -x^2 e^{-x} + 2 \left[ -x e^{-x} + e^{-x} \right] \right)_{0}^{\infty}
\]

Simplifying:

\[
= \frac{1}{2} \left[ 1 \right] + \frac{1}{2} \left[ 1 \right]
\]

\[
= 1 + 1 = 2.
\]

Thus, \( \text{Var}(X) \) for the given probability density function is 2.



\section*{2A The birthday problem}

The famous birthday problem is to find the smallest number of people one must ask to get an even chance that at least two people have the same birthday. To solve this, you can use the following steps. Find the probability that in a group of \( k \) people, no two have the same birthday. Let \( q \) be this probability. Then \( p = 1 - q \) is the probability that at least two people have the same birthday. Ignoring leap years, take the sample space \( S \) as all sequences of length \( k \) with each element being one of the 365 days in the year. Thus, there are \( 365^k \) elements in \( S \).

\begin{enumerate}
    \item[(a)] Find the total number of sequences with no common birthdays.
    
    \item[(b)] Assuming that each sequence is equally likely, show that
    \[
    q = \frac{365 \cdot 364 \cdot \cdots \cdot (365 - k + 1)}{365^k}
    \]

    \item[(c)] Write a computer program for calculating \( q \) for \( k = 2 \) to \( k = 50 \) and find the first \( k \) for which \( p \geq 0.5 \). This will give the least number of people we should ask to make it an even chance that at least two people will have the same birthday.
\end{enumerate}

\section*{Solution:}

\textbf{Solution:}

\textbf{(a)} Find the total number of sequences with no common birthdays.

Each of the \( k \) people can be assigned a birthday in the following way:

\begin{itemize}
    \item The first person has 365 choices (since there are 365 days in a year).
    \item The second person has 364 choices (since one day is already taken).
    \item The third person has 363 choices, and so on.
\end{itemize}

The total number of sequences where no two people have the same birthday is:

\[
365 \times 364 \times 363 \times \cdots \times (365 - k + 1).
\]

\textbf{(b)} Show that the probability \( q \) is given by the formula.

The total number of possible birthday assignments (without any restriction on shared birthdays) is \( 365^k \). Since there are 365 days for each person, the total number of ways to assign birthdays to \( k \) people is:

\[
365^k.
\]

Now, the probability \( q \) that no two people share the same birthday is the ratio of the number of sequences where no two people share the same birthday (from part (a)) to the total number of possible sequences:

\[
q = \frac{365 \times 364 \times 363 \times \cdots \times (365 - k + 1)}{365^k}.
\]

Thus, we have:

\[
q = \frac{365 \cdot 364 \cdot \cdots \cdot (365 - k + 1)}{365^k}.
\]

\textbf{(c)} Write a computer program for calculating \( q \) for \( k = 2 \) to \( k = 50 \) and find the first \( k \) for which \( p \geq 0.5 \).

To calculate \( q \) for \( k = 2 \) to \( k = 50 \), and find the smallest \( k \) for which \( p = 1 - q \geq 0.5 \), here is the pseudocode:

\begin{verbatim}
def birthday_problem():
    k = 0
    q = 1
    p = 0
    total_days = 365
    while p < 0.5:
        k += 1
        q *= (total_days - k + 1) / total_days
        p = 1 - q
        print(f"k = {k}, q = {q}, p = {p}")
        if p >= 0.5:
            return k

result = birthday_problem()
print(f"The smallest k for which p >= 0.5 is: {result}")
\end{verbatim}

By running this program, the smallest \( k \) for which \( p \geq 0.5 \) is typically \( k = 23 \). This means that in a group of 23 people, there is an even chance (50% or greater) that two people share the same birthday.

\begin{figure}[H] % The H option forces the figure to be placed exactly here
\centering
% Replace this with the correct path to your Table 1.9 image file
\includegraphics[width=0.9\textwidth]{D:/Masters NPU Semeter/3rd Semester/Statistics/Homework/Submit/Homework 2/figs/birthday.png}
\end{figure}

\end{document}